\documentclass[11pt]{article}
\usepackage[margin=0.88in]{geometry}
\usepackage[T1]{fontenc}
\usepackage{lmodern,amsmath,amssymb,amsthm,booktabs,array,longtable,microtype}
\usepackage[breaklinks=true,colorlinks=true,linkcolor=blue,urlcolor=blue,citecolor=blue]{hyperref}
\newtheorem{proposition}{Proposition}
\newcommand{\D}{\mathrm D}
\newcommand{\E}{\mathbb E}
\newcommand{\Var}{\operatorname{Var}}
\newcommand{\tr}{\operatorname{tr}}
\newcommand{\N}{\mathcal N}
\title{A proposal to repair the M13 particle-flow likelihood and score\\
while preserving LGSSM and KSC-SV}
\author{BayesFilter working proposal}
\date{4 October 2026}
\begin{document}
\maketitle

The proposed repair is to sum over possible ancestors in the importance
correction, while retaining the existing UKF covariance lifecycle, LEDH flow,
Contract E reset, and analytical differentiation. This is the marginal
particle-filter construction applied to the actual affine flow proposals.
It changes the finite likelihood approximation and therefore needs a distinct
route identity and its own tuning. It is not the unchanged Li--Coates
Algorithm 1.

The purpose remains testing \emph{surrogate HMC}. A particle-filter score may
supply a force for proposals, but the Metropolis calculation must use the
declared endpoint target. A better filter score can make those proposals more
useful; it does not by itself establish posterior correctness or convergence.

The mathematics below establishes the proposed one-step integration identity
and the total derivative of the resulting finite program. It also establishes
why a universal finite-sample improvement claim would be wrong. The existing
LGSSM and KSC-SV computations are preserved by retaining their existing route,
then checking that preservation through the actual consumers. The new M13
route must earn admission by measurement.

\section{The observed failure and the quantity to repair}

The October 4 localization traced the M13 consumer through the fused batch
route into the canonical scalar value and analytical-score implementation.
At all eight diagnosed parameter points, the diagnostic trace and full
consumer agreed in FP64. The analytical score also agreed with finite
differences of its \emph{own} implemented likelihood to about
$2.6\,10^{-6}$ in the preceding score check. This rules out a large missing
derivative term as the explanation supported by those checks. It does not
establish accuracy relative to the statistical model.

M13 has a four-dimensional state and two observations (range and bearing).
At the diagnosed low-process-noise point, the first observation already shows
the failure:
\begin{center}
\begin{tabular}{lrrr}
\toprule
Method & Log likelihood & First score coordinate & Particle ESS\\
\midrule
Refined independent quadrature & $-0.13791267$ & $0.01149807$ & ---\\
Same-design bootstrap & $-0.13013616$ & $0.00812633$ & $205.04$\\
Current LEDH & $-2.22116858$ & $3.57681741$ & $6.66$\\
\bottomrule
\end{tabular}
\end{center}
There are 4,096 particles. The second score coordinates are respectively
$-0.10229746$, $-0.10443582$, and $0.19335105$.
The diagnostic parameter is
$\theta=(-0.73272085,-0.01465508)$ in the recorded M13 coordinates.

Before applying the observation density, the importance estimate of an
integral equal to one is only $0.00344047$. The first UKF position variance
is about 116.74 times the conditional process variance. The root mean square
flow displacement, measured in process-covariance Mahalanobis units, is
20.61. These observations support a specific mechanism: a child can be moved
far from its own ancestor's narrow transition density. Its assigned
ancestor then explains it poorly, even when other ancestors could explain
the same location.

The discrepancy occurs before the first reset. Preserving the computed
cloud's mean and covariance during reset cannot reconstruct filtering
probability that the importance approximation has already lost.
Replacing the UKF covariance with the process covariance is not justified
by this observation: those covariances describe different conditional
distributions, and the UKF lifecycle is part of the specified algorithm.

The quadrature was refined in both grid size and finite-difference step.
This is a numerical reference, not an exact closed-form M13 score oracle.
The eight points are now development data. Their results are descriptive
and cannot serve as untouched promotion evidence.

\section{Specify the finite mixture before changing its weights}

Fix one time step and condition on the complete previous cloud and all
offline numerical controls. Write the ancestors and their normalized
positive weights as $(a_j,\omega_j)_{j=1}^N$, where
$\sum_j\omega_j=1$. For M13,
\begin{equation}
 f_j(x)=p_\theta(x\mid a_j)=\N(x;m_j,Q_\theta),\qquad
 p_N(x)=\sum_{j=1}^N\omega_j f_j(x),\qquad
 Z_N=\int g_\theta(x)p_N(x)\,dx .
 \label{eq:predictive-mixture}
\end{equation}
Here $m_j=Aa_j$ with the model's declared transition, and $g_\theta(x)$ is
the implemented range--bearing observation density, including its angle
convention. $Z_N$ is the normalizer of this \emph{finite predictive
mixture}; it is not automatically the exact model's predictive likelihood.

In the current M13 implementation, each auxiliary flow starts at the
zero-process-noise prediction. Conditional on the previous cloud, its
coefficients do not depend on the current process innovation. The composed
Euler map acting on that innovation's child is therefore affine:
\begin{equation}
 F_j(u)=B_j u+c_j,\qquad
 q_j(x)=\frac{f_j\!\left(B_j^{-1}(x-c_j)\right)}
                   {|\det B_j|}.
 \label{eq:affine-density}
\end{equation}
Assume $Q_\theta$ is positive definite and $B_j$ nonsingular.
These conditions make $q_j$ a normalized pushforward density.
Nonsingularity must hold in the neighborhood in which a score is claimed.
For a general nonlinear map, a nonsingular Jacobian at the sampled point
alone would not establish global invertibility; the conditional affine
structure is doing real work here.

This construction uses a full-dimensional density with respect to Lebesgue
measure. A structurally singular transition, such as a deterministic
coordinate in AR(2), needs a separate innovation-space measure derivation.
The proposed full-rank route must reject that unsupported case, rather
than add an undeclared ridge or call it the same density.

With one child $X_i=F_i(U_i)$ from each ancestor, $U_i\sim f_i$, the current
correction has the form
\begin{equation}
 V_i^{\rm old}=\omega_i g_\theta(X_i)
                         \frac{f_i(X_i)}{q_i(X_i)} .
 \label{eq:old-weight}
\end{equation}
Define the marginal proposal and proposed weights by
\begin{equation}
 q_N(x)=\sum_j\omega_j q_j(x),\qquad
 V_i^{\rm mar}=\omega_i g_\theta(X_i)
                       \frac{p_N(X_i)}{q_N(X_i)},\qquad
 \widehat Z_{\rm mar}=\sum_i V_i^{\rm mar}.
 \label{eq:marginal-weight}
\end{equation}
Both the numerator and denominator change to mixtures. Replacing only one
of them computes a different correction.

The leading $\omega_i$ is correct for one draw from each component and the
weighted sum displayed here. If instead $M$ labels are drawn independently
from $\omega$, the estimator is the average of $M$ marginal weights, with
leading factor $1/M$, not a second factor $\omega_{J_i}$.
This distinction must survive any implementation refactor.

\section{What marginalization proves, and what it does not}

\begin{proposition}[Conditional integration identity]
\label{prop:integration}
Suppose each $X_i$ has conditional marginal density $q_i$, $g_\theta p_N$
is integrable, and $q_N>0$ wherever $g_\theta p_N>0$. Then
$\E[\widehat Z_{\rm mar}\mid (a,\omega)]=Z_N$.
\end{proposition}
\begin{proof}
Linearity of expectation and the definition of the mixture give
\begin{align}
 \E\widehat Z_{\rm mar}
 &=\sum_i\omega_i\int q_i(x)
                    \frac{g_\theta(x)p_N(x)}{q_N(x)}\,dx \notag\\
 &=\int\frac{\sum_i\omega_i q_i(x)}{q_N(x)}
                        g_\theta(x)p_N(x)\,dx
 =Z_N .
 \label{eq:integration-proof}
\end{align}
Independence between children is unnecessary for this expectation identity.
\end{proof}

Setting $g_\theta=1$ gives the pre-observation normalization target one.
This provides a cheap diagnostic of integration error. The statement is
conditional on the specified cloud and sampling law. A frozen Halton
trajectory is not a random experiment satisfying those assumptions merely
because its inputs once came from a randomized construction. Neither this
proof nor the marginal-particle-filter literature licenses a claim that the
log likelihood, its score, or the complete likelihood product after the
repository's deterministic Contract E and GenUT operations is unbiased.
Unbiasedness of a positive likelihood estimator would not itself imply
unbiasedness after taking its logarithm.

\begin{proposition}[Variance reduction in the random-label experiment]
\label{prop:variance}
Draw $J$ with probabilities $\omega_j$, then $X\mid J=j\sim q_j$.
Let $W=g_\theta(X)f_J(X)/q_J(X)$ and
$W_m=g_\theta(X)p_N(X)/q_N(X)$. If the relevant second moments are finite,
$\Var(W_m)\leq\Var(W)$.
\end{proposition}
\begin{proof}
Bayes' rule gives
$\Pr(J=j\mid X=x)=\omega_j q_j(x)/q_N(x)$. Hence
\begin{equation}
 \E[W\mid X=x]
 =\sum_j\frac{\omega_j q_j(x)}{q_N(x)}
                    \frac{g_\theta(x)f_j(x)}{q_j(x)}
 =\frac{g_\theta(x)p_N(x)}{q_N(x)}=W_m(x).
 \label{eq:rb-proof}
\end{equation}
Applying the law of total variance,
\begin{equation}
 \Var(W)-\Var(W_m)=\E[\Var(W\mid X)]\geq0.
 \label{eq:variance-decomposition}
\end{equation}
For two components, with $0<w<1$ and $q_1,q_2>0$, the pointwise
second-moment difference has the explicit nonnegative form
\begin{equation}
 \frac{wf_1^2}{q_1}+\frac{(1-w)f_2^2}{q_2}
 -\frac{(wf_1+(1-w)f_2)^2}{wq_1+(1-w)q_2}
 =
 \frac{w(1-w)(f_1q_2-f_2q_1)^2}
      {q_1q_2(wq_1+(1-w)q_2)} .
 \label{eq:two-component-square}
\end{equation}
Multiply by $g_\theta^2$ and integrate to obtain the corresponding
second-moment comparison.
\end{proof}

This is the Rao--Blackwell argument behind the marginal particle filter
\cite[Sec.~3, Fig.~3, Sec.~3.1 and Prop.~1]{klaas}.
Its sampling experiment differs from a stratified allocation with exactly
one child per component. The distinction prevents an invalid guarantee.

For an explicit counterexample, consider two states $A,B$ with
\[
 g=(1,2),\quad f_1=(9/10,1/10),\quad f_2=(1/10,9/10),
 \quad \omega_1=\omega_2=1/2,
\]
and choose $q_1=(9/11,2/11)$, $q_2=(1/19,18/19)$.
Then $g f_1/q_1=11/10$ and $g f_2/q_2=19/10$, constant within
each component. With independent draws, one from each component,
the original estimate is exactly $3/2$ and has zero variance.
But $p_N=(1/2,1/2)$ and $q_N=(91/209,118/209)$, so the marginal
weights at $A,B$ are $209/182$ and $209/118$. Direct calculation gives
\begin{align}
 \E\widehat Z_{\rm mar}&=3/2,\notag\\
 \Var(\widehat Z_{\rm mar})
 &=\frac14\left(\frac9{11}\frac2{11}
                    +\frac1{19}\frac{18}{19}\right)
            \left(\frac{209}{182}-\frac{209}{118}\right)^2
 =\frac{555264}{28826161}>0 .
 \label{eq:counterexample}
\end{align}
Thus the proposed marginal weights can be worse under stratification.
There is no universal finite-design non-deterioration theorem to transfer
to LGSSM, KSC-SV, or M13. Even the valid random-label theorem says nothing
by itself about the variance or accuracy of log-likelihood derivatives.

For M13, the reason to test the repair is narrower: the diagnosed
ancestor-conditional ratios are extremely poorly integrated, whereas
a mixture can explain a transported particle using other ancestors.
Whether that opportunity is sufficient is an empirical question.

\section{The score must differentiate the same finite program}

Let $\D$ denote the total derivative in an arbitrary parameter direction,
holding the base design and all offline controls fixed. It includes the
dependence of previous particles, weights, current children, UKF moments,
flow coefficients and process/observation covariances on the parameter.
Work on a neighborhood where these functions are differentiable, densities
are positive and the declared numerical status is valid. An angle-branch
boundary or an invalid factorization requires an explicit status; it cannot
be hidden in a smooth-score assertion.

At a child location $x=X_i(\theta)$ define the responsibilities
\[
 r_j^f=\frac{\omega_j f_j(x)}{p_N(x)},\qquad
 r_j^q=\frac{\omega_j q_j(x)}{q_N(x)}.
\]
Differentiating a finite sum and dividing by that sum yields
\begin{align}
 \D\log p_N(x)
   &=\sum_j r_j^f\{\D\log\omega_j+\D\log f_j(x)\},\notag\\
 \D\log q_N(x)
   &=\sum_j r_j^q\{\D\log\omega_j+\D\log q_j(x)\}.
 \label{eq:mixture-score}
\end{align}
For example,
$\D p_N=\sum_j\{(\D\omega_j)f_j+\omega_j\D f_j\}$;
division by $p_N$ gives the first line. Responsibilities appear as
coefficients in this first derivative; no extra derivative of the
responsibility is added. That extra term would belong to a second derivative.

The proposed log-weight derivative is therefore
\begin{equation}
 \D\log V_i^{\rm mar}
 =\D\log\omega_i+\D\log g_\theta(X_i)
       +\D\log p_N(X_i)-\D\log q_N(X_i).
 \label{eq:weight-score}
\end{equation}
With $\bar V_{ti}=V_{ti}/\sum_k V_{tk}$, the finite log likelihood
and its derivative are
\begin{equation}
 \widehat\ell(\theta)=\sum_t\log\!\sum_i V_{ti}(\theta),\qquad
 \D\widehat\ell(\theta)=
             \sum_t\sum_i\bar V_{ti}\D\log V_{ti}.
 \label{eq:finite-score}
\end{equation}
This is an ordinary finite-program chain rule, not differentiation under
a probability integral. It proves which derivative an implementation must
compute, rather than accuracy relative to the model likelihood.

To evaluate a proposal component at an arbitrary child, put
$v_j=B_j^{-1}(x-c_j)$. Differentiating $B_jv_j=x-c_j$ gives
\begin{equation}
 \D v_j=B_j^{-1}\{\D x-\D c_j-(\D B_j)v_j\}.
 \label{eq:inverse-tangent}
\end{equation}
On a nonsingular neighborhood the determinant has constant sign, and
Jacobi's identity gives
\begin{equation}
 \D\log q_j(x)=\D\log f_j(v_j)
                      -\tr(B_j^{-1}\D B_j).
 \label{eq:proposal-score}
\end{equation}
In particular, omitting the inverse-map tangent is wrong relative to the
total-score claim.

For a Gaussian component let $e=x-m$ and $u=Q^{-1}e$.
From $\log f=-\tfrac12\log\det Q-\tfrac12 e^\top Q^{-1}e+\mathrm{const}$,
$\D Q^{-1}=-Q^{-1}(\D Q)Q^{-1}$ and
$\D\log\det Q=\tr(Q^{-1}\D Q)$, obtain
\begin{equation}
 \D\log f(x)
 =u^\top(\D m-\D x)
  +\tfrac12u^\top(\D Q)u-\tfrac12\tr(Q^{-1}\D Q).
 \label{eq:gaussian-score}
\end{equation}
The terms in $\D x$, $\D m$ and $\D Q$ all matter. In the denominator,
replace $x$ by $v_j$ and use \eqref{eq:inverse-tangent}.

The composed maps and their derivatives can be accumulated in the existing
flow loop. For Euler step $k$ let $M_k=I+\epsilon_k A_k$, with the step
schedule fixed offline. Then
\begin{align}
 B_{k+1}&=M_k B_k,& c_{k+1}&=M_k c_k+\epsilon_k b_k,
 \label{eq:map-composition}\\
 \D B_{k+1}&=(\D M_k)B_k+M_k\D B_k,&
 \D c_{k+1}&=(\D M_k)c_k+M_k\D c_k+\epsilon_k\D b_k .
 \label{eq:map-tangent}
\end{align}
Initialize $B_0=I$, $c_0=0$ and both tangents at zero.
The existing analytical $A_k,b_k,\D A_k,\D b_k$ must be reused.
Do not construct a simpler map in the batch lane or drop the UKF/auxiliary
dependence. Matrix multiplication order must be preserved.

After weighting, the existing reset and correction still have to propagate
the total derivative. In particular Contract E requires both the direct
moment/weight terms and the terms through transport; changing the weights
does not remove either dependence. The derivation is inductive: if the
incoming cloud and UKF tangents are correct, the displayed formulas give
correct weight tangents; the existing reset/correction must then return
correct outgoing tangents. Wiring and end-to-end finite differences must
check this induction. No new implementation exists yet, so that call-chain
claim is not checked for the proposed route.

\section{An optional defensive component if marginalization is insufficient}

A mixture over flowed ancestors can still miss important regions. A second
candidate, only if that problem persists, uses the actual proposal
\begin{equation}
 r_j(x)=\alpha f_j(x)+(1-\alpha)q_j(x),\qquad
 r_N(x)=\alpha p_N(x)+(1-\alpha)q_N(x),\quad 0<\alpha\leq1,
 \qquad W_\alpha(x)=\frac{g_\theta(x)p_N(x)}{r_N(x)} .
 \label{eq:defensive}
\end{equation}
Samples must be drawn from $r_j$, for example by an offline frozen,
parameter-independent Bernoulli label and the corresponding base
innovation. Merely changing the denominator while continuing to sample
only from $q_j$ is wrong. As with the earlier identity, the randomized
sampling-law claim and any particular frozen realization are distinct.
This proposal keeps $\alpha$ fixed; parameter-dependent or online
selection would require another method and derivative derivation.

Since $r_N\geq\alpha p_N$ pointwise, wherever $p_N>0$,
\begin{equation}
 \frac{p_N}{r_N}\leq\frac1\alpha,\qquad
 \E_{r_N}W_\alpha^2
 =\int\frac{g_\theta^2 p_N^2}{r_N}
 \leq\frac1\alpha\E_{p_N}g_\theta^2 .
 \label{eq:defensive-bound}
\end{equation}
Provided $\E_{p_N}g_\theta^2<\infty$, subtraction of $Z_N^2$ yields
\begin{equation}
 \Var_{r_N}(W_\alpha)
 \leq\frac1\alpha\Var_{p_N}(g_\theta)
           +\left(\frac1\alpha-1\right)Z_N^2 .
 \label{eq:defensive-variance}
\end{equation}
The last term must not be dropped. The bound is on a random marginal draw,
not a uniform guarantee about stratified quadrature error, score error,
or the adequacy of the previous cloud. It does not guarantee high ESS.
At $\alpha=1$ this is bootstrap sampling; at $\alpha=0$ one obtains the
unprotected marginal proposal, outside the displayed bound.

This is the basic defensive-importance argument discussed by
Owen and Zhou \cite[Sec.~3, especially Eq.~(10), and Sec.~3.1]{owen}.
Their stronger results use additional control variates; those are not
part of this proposal and their guarantees must not be borrowed.
A defensive mixture can deteriorate an already excellent importance
proposal. No numerical value of $\alpha$ has been justified here.
Its choice would need scope-specific offline calibration and a separate
non-harm assessment, never selection inside HMC.

\section{Preserve the original LGSSM and KSC-SV values and scores}

Preservation has a precise and achievable meaning. Existing qualified
scopes keep their existing algorithm, prepared observations, base designs,
settings, coordinate transforms and consumer endpoints. The proposed
correction is an explicit optional mode selected before tracing. It must
not be a parameter-dependent switch and must not silently become the
factory default. The mode is generic; this does not justify a reduced,
model-specific implementation of the canonical flow.

\begin{proposition}[Preservation by unchanged execution]
\label{prop:preserve}
Let $A_{\rm old}(\theta;\xi,c)$ denote the current finite program with fixed
inputs $\xi$ and controls $c$. If the new factory's legacy mode calls
that same program on those same inputs, with no changed operations,
then on its valid differentiable domain
\begin{equation}
 A_{\rm legacy}(\theta;\xi,c)=A_{\rm old}(\theta;\xi,c),\qquad
 \D A_{\rm legacy}(\theta;\xi,c)=\D A_{\rm old}(\theta;\xi,c).
 \label{eq:preserve}
\end{equation}
\end{proposition}
\begin{proof}
The function computed in legacy mode is the same function. Equal
differentiable functions on a neighborhood have equal derivatives there.
If the implementation returns a separate analytical score, preserving its
operations and inputs also preserves that returned score, including any
pre-existing error. This proposition does not newly certify that score.
\end{proof}

The implementation obligation is substantive: a refactor can change
reduction order, random-number consumption, precision or a shared helper.
The proof cannot certify those details before code exists. Freeze the
current \emph{post-fix} source closure and actual outputs, then exercise the
real consumers to verify them. Bitwise parity is the goal where execution
is deterministic. Any nondeterministic FP32 tolerance must come from
measured same-build repeatability, not from a convenient new error allowance.

The original scalar LG1 model has state and observation dimensions one:
$x_t=\phi x_{t-1}+\epsilon_t$ with process variance one, observation
variance $1/2$, and stationary initial variance $(1-\phi^2)^{-1}$.
Its exact Kalman value and analytical derivative are the independent
reference. Preserve the derivative of the initial variance,
$2\phi/(1-\phi^2)^2$, and the existing bounded-parameter chart and prior.
The regression inventory must include every already qualified LGSSM
consumer/scope, not infer coverage of all linear models from this one case.

The protected KSC consumer is the seven-component approximation to the
log-squared-return observation law:
\begin{equation}
 g_\theta(z\mid x)=
       \sum_{k=1}^7\pi_k\,\N(z;x+2\log\beta+\mu_k,v_k).
 \label{eq:ksc-law}
\end{equation}
It is distinct from the raw-return Gaussian stochastic-volatility
observation density. For residual
$e_k=z-x-2\log\beta-\mu_k$ and fixed mixture constants, each component has
\begin{equation}
 \D\log g_k=\frac{e_k}{v_k}
                  \left(\D x+\frac{2\,\D\beta}{\beta}\right).
 \label{eq:ksc-tangent}
\end{equation}
The recently repaired $2\,\D\beta/\beta$ term must remain.
The score of \eqref{eq:ksc-law} is the responsibility-weighted sum of these
component scores. The $T=2$ KSC fixture has an independent 49-component
Gaussian likelihood after integrating its latent states; it is exact for
this declared KSC mixture model, not for the unapproximated SV model.
Preserve the stationary state variance and its derivative, chart,
prior and observation mixture.

For both protected families, compare separately the physical likelihood,
physical score, transformed target and transformed score. Check FP64
against the frozen post-fix implementation and the independent reference.
Check the actual FP32 GPU/XLA consumer using the repository memory policy.
Require the previous error envelope against the reference to be maintained
for both values and scores; an improved value cannot excuse a worse score.
Statuses, finiteness and deterministic outputs are part of preservation.
Repeated builds or randomizations need paired uncertainty where relevant.

This proposal does not assert that marginal weighting will be as accurate
as the old computation on every LGSSM or KSC scope. Migration would require
new evidence. Retaining the original computation is what guarantees that
testing the M13 candidate does not degrade those working results.

\section{Implementation and the smallest discriminating experiments}

The current M13 call chain is
\path{ledh_range_bearing_target_tf.py}
$\to$ \path{ledh_canonical_batch_fused_tf.py}
$\to$ \path{ledh_canonical_score_tf.py}.
In the last module the current conditional density is formed near line
475, the flow near 594, and the importance correction near 611--652.
The affine flow updates are near 1167--1403. Line numbers refer to the
current dirty checkout and will move if it changes.
The previous executable trace parity establishes the baseline chain,
not the unimplemented extension.

The proposed implementation would return $B,c,\D B,\D c$ from the shared
flow and evaluate the two mixture densities and their tangents at all
children. Use log-sum-exp and responsibility reductions with stable
TensorFlow graphs; keep the analytical score as the claim-bearing path.
Autodiff and finite differences are independent diagnostics only.
Tiled sums must equal the dense mathematical sums. At $N=4096$ the
repository chunk policy gives $K=2048$. There are $N^2$ component/child
pairs per step: tiling controls memory, not this arithmetic cost.
A truncation to nearby ancestors would be an additional approximation,
not the exact method derived here.

The following sequence belongs to the existing master's M13 repair phase.
It is proposed work, not an already executed new sampling campaign.

\begin{enumerate}
\item Freeze the current post-fix LGSSM/KSC source, observations, controls
and output baselines. Add an optional generic correction mode. Run
consumer-level dispatch and value/score preservation checks before
enabling the M13 candidate. Preserve canonical identity through the
repository-owned factory; a changed source closure must be recorded,
not hidden behind a caller-stamped identifier.

\item Check affine composition, inverse densities and analytical
tangents on small full-rank cases, including noncommuting matrices.
Check batch/scalar and tiled/dense parity and invalid-map status handling.
Test the complete finite likelihood against directional finite differences
on stable branches. This establishes engineering correctness only.

\item At $T=1$, compare original LEDH, marginal LEDH, the same-design
bootstrap and refined quadrature. Use the diagnosed low-$Q$ case first,
then all eight development points, reporting low noise, center and tails
separately. Record the pre-observation normalizer, weight concentration,
value error and both score errors. Movement of normalization toward one
is explanatory; simultaneous value/score agreement with a qualified
reference is the local accuracy criterion.

\item If proposal coverage still fails, evaluate the defensive component
as a separately tuned candidate, including $\alpha=1$ bootstrap and
$\alpha=0$ marginalization as comparators. Choose no online thresholds
or defensive proportions. A failed candidate triggers the next
discriminating repair, not automatic rejection of particle flow.

\item Before a $T=20$ accuracy decision, qualify the independent reference
\emph{score} by refinement and independent-design uncertainty. Stability
of its likelihood alone is insufficient. Use fresh calibration/validation
partitions for tuning and untouched cases for any promotion. Do not select
settings using the eight now-inspected points as a final holdout.

\item Only then test the actual surrogate-HMC consumer with its declared
endpoint, the supported public HMC tuner, and the repository's adaptive
warm-up and cumulative sampling procedure. Do not introduce another
arbitrary fixed terminal sample count. Posterior agreement and declared
sampler diagnostics remain the campaign's substantive criteria.
\end{enumerate}

The correct tuner must be chosen using
\path{docs/reference/hmc-tuning-interface.md} and
\path{HMC_TUNING_INTERFACE_CAPABILITIES}.
The existing neural-force runner binding supports the public
\path{tune_hmc_kernel} route; a stage helper is not a tuner, and an
arbitrary force must not be routed through the fixed-transport tuner.

Every new scope still owes offline LEDH tuning. FP32 GPU/XLA remains the
default execution target, with memory growth verified before device
initialization; FP64 is a reference arm. No new tolerance is selected from
the observed failures. For HMC materiality, useful explanatory quantities
are local log-likelihood-ratio error
$\delta(\theta+h)-\delta(\theta)$ and score error in the declared HMC
coordinates, such as $\|L^\top(\widehat s-s_{\rm ref})\|$ for
$\theta=\theta_0+Lz$. Neither substitutes for downstream sampling evidence.

\section{Mathematical audit and limits of this proposal}

Manual proofs above establish the stated conditional identities. MathDevMCP
has verified the two-component variance identity, conditional-expectation
identity, rational mixture derivative, Gaussian scalar tangent,
inverse-map scalar tangent, KSC tangent and exact counterexample arithmetic
using its symbolic backend. The separate audit record preserves exact
inputs, outputs, parser limitations and document-audit results.

An attempted nested differentiation expression was not encoded correctly by
the tool. It was replaced by an explicit rational derivative identity.
A separate high-level request to refute universal stratified
non-deterioration returned an inconclusive wrapper status despite an
exact nonzero backend difference. The explicit counterexample and its
checked rational arithmetic carry the argument; the wrapper abstention
is not represented as a proof certificate.

These checks do not certify the unimplemented TensorFlow call chain,
empirical improvement, universal non-deterioration, a full-horizon exact
M13 score or HMC convergence. This stage changes no algorithmic source.
The review question is whether this is a mathematically valid, honestly
bounded and implementable M13 repair proposal with an adequate preservation
protocol for the working LGSSM and KSC-SV consumers.

\begin{thebibliography}{9}
\bibitem{li}
Y. Li and M. Coates (2017).
Particle Filtering with Invertible Particle Flow.
\emph{IEEE Transactions on Signal Processing}, 65(16), 4102--4116.
Algorithm 1, Eqs.~(19)--(20), Lemma IV.2, Theorem IV.3 and Lemma IV.4.
\url{https://arxiv.org/abs/1607.08799}.
The local implementation was inspected; the original-author software
link was unavailable during this audit. No author-code equivalence
claim follows from the paper inspection.

\bibitem{klaas}
M. Klaas, N. de Freitas and A. Doucet (2005).
Toward Practical $N^2$ Monte Carlo: the Marginal Particle Filter.
\emph{Proceedings of UAI}, 308--315.
Sec.~3 and Fig.~3 (mixture weights); Sec.~3.1 and Proposition 1
(random-label variance argument).
\url{https://www.cs.ubc.ca/~arnaud/klass_defreitas_doucet_marginalparticlefilterUAI2005.pdf}.

\bibitem{owen}
A. B. Owen and Y. Zhou (2000).
Safe and Effective Importance Sampling.
\emph{Journal of the American Statistical Association}, 95(449), 135--143.
Sec.~3, Eq.~(10) and Sec.~3.1.
\url{https://artowen.su.domains/reports/seis.pdf}.
\end{thebibliography}
\end{document}
